% --- Archon physics formalization source begin ---
% archon:physics
% archon:covers PhyXMiniProblems/problem_phyx_mini_0903.lean
% archon:source-report reports/phyx_mini/problem_phyx_mini_0903.source.json

\chapter{Physics problem phyx\_mini\_0903}
\label{ch:PhyXMiniProblems_problem_phyx_mini_0903}

\paragraph{Problem source.}
\#\# Question

What is the magnitude of the electric fields at point 1?

\#\# Answer choices

- A: A: 1.25 \textbackslash{}times 10\textasciicircum{}\{3\}N/C

- B: B: 1.35 \textbackslash{}times 10\textasciicircum{}\{5\}N/C

- C: C: 1.8 \textbackslash{}times 10\textasciicircum{}\{5\}N/C

- D: D: 5.33 \textbackslash{}times 10\textasciicircum{}\{5\}N/C

\#\# Figure caption (auxiliary; use the image as primary evidence)

The image depicts a geometric arrangement of three points forming a triangle. Here's a detailed description:

1. **Objects and Layout:**
   - There are three distinct points forming a triangle.
   - The top point is marked by a teal circle with a negative sign inside, labeled as "-2.0 nC" indicating a negative charge of 2.0 nanocoulombs.
   - The bottom two points are black dots labeled as "1" and "2."

2. **Distances and Connections:**
   - The triangle is equilateral with each side measuring 1.0 cm. 
   - The lines between the points are dashed, indicating the connections or distances between them.

3. **Text:**
   - The measurements "1.0 cm" are labeled on each side of the triangle.
   - The charge "-2.0 nC" is noted above the teal circle.

The triangle outlines the relationships among the charged point at the top and the bottom two points labeled "1" and "2."

\#\# Dataset metadata

Electrostatics; Spatial Relation Reasoning, Physical Model Grounding Reasoning

\paragraph{Recorded answer/context.}
C: C: 1.8 \textbackslash{}times 10\textasciicircum{}\{5\}N/C

\paragraph{Figure/image path.}
/root/proposal\_for\_physic/science-mango/phyx\_mini\_run/phyx\_data/test\_image/903.png

\paragraph{Formalization target.}
create a compiling Lean file with sorry bodies at `PhyXMiniProblems/problem\_phyx\_mini\_0903.lean`. The Lean declarations must preserve the physical quantities, dimensions or dimensional roles, figure labels, governing-law hypotheses, and final relation expressed by this problem.
Use Mathlib/Physlib names found through LeanExplore where available. If a physics API is missing, introduce faithful local abstractions rather than scalar placeholder aliases.

\begin{theorem}[Physics formalization target]
\label{thm:physics:phyx_mini_0903:target}
\lean{PhyXMiniProblems.ProblemPhyXMini0903.problem_phyx_mini_0903}
\uses{def:physics:phyx-mini-0903:phyxminiproblems-problemphyxmini0903-electricfieldmagnitudeinnewtonspercoulomb, def:physics:phyx-mini-0903:phyxminiproblems-problemphyxmini0903-pointchargetrianglesetup, def:physics:phyx-mini-0903:phyxminiproblems-problemphyxmini0903-matchessinglepointchargescenario, def:physics:phyx-mini-0903:phyxminiproblems-problemphyxmini0903-matchessuppliedelectricfieldtrianglefigure, def:physics:phyx-mini-0903:phyxminiproblems-problemphyxmini0903-usesschoolcoulombconstant, def:physics:phyx-mini-0903:phyxminiproblems-problemphyxmini0903-hasphysicalpointchargetriangleparameters, def:physics:phyx-mini-0903:phyxminiproblems-problemphyxmini0903-satisfiesplanarpointchargecoulomblaw}
The assigned autoformalize agent should translate this physics problem into Lean declarations in the covered file, with theorem and lemma proof bodies written as `by sorry`.
\end{theorem}
\begin{proof}
This is an autoformalization task, not a proof task. Produce statements that can later be proved by the physics prover without weakening the physical model.
\end{proof}

% --- Archon declaration topology begin ---
\section{Declaration topology}
\begin{definition}[electric Field Dimension]
  \label{def:physics:phyx-mini-0903:phyxminiproblems-problemphyxmini0903-electricfielddimension}
  \lean{PhyXMiniProblems.ProblemPhyXMini0903.electricFieldDimension}
  The physical dimension M L T⁻² C⁻¹ of an electric field
\end{definition}
\begin{proof}
  This is the declared model definition of electric Field Dimension.
\end{proof}
\begin{definition}[Plane Vector]
  \label{def:physics:phyx-mini-0903:phyxminiproblems-problemphyxmini0903-planevector}
  \lean{PhyXMiniProblems.ProblemPhyXMini0903.PlaneVector}
  Vectors in the two-dimensional plane of the supplied diagram
\end{definition}
\begin{proof}
  This is the declared model definition of Plane Vector.
\end{proof}
\begin{definition}[Signed Charge Quantity]
  \label{def:physics:phyx-mini-0903:phyxminiproblems-problemphyxmini0903-signedchargequantity}
  \lean{PhyXMiniProblems.ProblemPhyXMini0903.SignedChargeQuantity}
  A signed, unit-independent electric charge
\end{definition}
\begin{proof}
  This is the declared model definition of Signed Charge Quantity.
\end{proof}
\begin{definition}[Plane Position Quantity]
  \label{def:physics:phyx-mini-0903:phyxminiproblems-problemphyxmini0903-planepositionquantity}
  \lean{PhyXMiniProblems.ProblemPhyXMini0903.PlanePositionQuantity}
  \uses{def:physics:phyx-mini-0903:phyxminiproblems-problemphyxmini0903-planevector}
  A unit-independent planar position, expressed relative to a chosen origin
\end{definition}
\begin{proof}
  This is the declared model definition of Plane Position Quantity.
\end{proof}
\begin{definition}[Length Quantity]
  \label{def:physics:phyx-mini-0903:phyxminiproblems-problemphyxmini0903-lengthquantity}
  \lean{PhyXMiniProblems.ProblemPhyXMini0903.LengthQuantity}
  A nonnegative, unit-independent physical length
\end{definition}
\begin{proof}
  This is the declared model definition of Length Quantity.
\end{proof}
\begin{definition}[Plane Electric Field Quantity]
  \label{def:physics:phyx-mini-0903:phyxminiproblems-problemphyxmini0903-planeelectricfieldquantity}
  \lean{PhyXMiniProblems.ProblemPhyXMini0903.PlaneElectricFieldQuantity}
  \uses{def:physics:phyx-mini-0903:phyxminiproblems-problemphyxmini0903-electricfielddimension, def:physics:phyx-mini-0903:phyxminiproblems-problemphyxmini0903-planevector}
  A unit-independent planar electric-field vector
\end{definition}
\begin{proof}
  This is the declared model definition of Plane Electric Field Quantity.
\end{proof}
\begin{definition}[charge In Coulombs]
  \label{def:physics:phyx-mini-0903:phyxminiproblems-problemphyxmini0903-chargeincoulombs}
  \lean{PhyXMiniProblems.ProblemPhyXMini0903.chargeInCoulombs}
  \uses{def:physics:phyx-mini-0903:phyxminiproblems-problemphyxmini0903-signedchargequantity}
  Coherent-SI readout of a signed charge in coulombs
\end{definition}
\begin{proof}
  This is the declared model definition of charge In Coulombs.
\end{proof}
\begin{definition}[charge In Nanocoulombs]
  \label{def:physics:phyx-mini-0903:phyxminiproblems-problemphyxmini0903-chargeinnanocoulombs}
  \lean{PhyXMiniProblems.ProblemPhyXMini0903.chargeInNanocoulombs}
  \uses{def:physics:phyx-mini-0903:phyxminiproblems-problemphyxmini0903-signedchargequantity, def:physics:phyx-mini-0903:phyxminiproblems-problemphyxmini0903-chargeincoulombs}
  Nanocoulomb readout used by the charge label in the image
\end{definition}
\begin{proof}
  This is the declared model definition of charge In Nanocoulombs.
\end{proof}
\begin{definition}[position In Meters]
  \label{def:physics:phyx-mini-0903:phyxminiproblems-problemphyxmini0903-positioninmeters}
  \lean{PhyXMiniProblems.ProblemPhyXMini0903.positionInMeters}
  \uses{def:physics:phyx-mini-0903:phyxminiproblems-problemphyxmini0903-planevector, def:physics:phyx-mini-0903:phyxminiproblems-problemphyxmini0903-planepositionquantity}
  Coherent-SI coordinate vector of a planar position, in metres
\end{definition}
\begin{proof}
  This is the declared model definition of position In Meters.
\end{proof}
\begin{definition}[length In Meters]
  \label{def:physics:phyx-mini-0903:phyxminiproblems-problemphyxmini0903-lengthinmeters}
  \lean{PhyXMiniProblems.ProblemPhyXMini0903.lengthInMeters}
  \uses{def:physics:phyx-mini-0903:phyxminiproblems-problemphyxmini0903-lengthquantity}
  Coherent-SI readout of a nonnegative length, in metres
\end{definition}
\begin{proof}
  This is the declared model definition of length In Meters.
\end{proof}
\begin{definition}[length In Centimeters]
  \label{def:physics:phyx-mini-0903:phyxminiproblems-problemphyxmini0903-lengthincentimeters}
  \lean{PhyXMiniProblems.ProblemPhyXMini0903.lengthInCentimeters}
  \uses{def:physics:phyx-mini-0903:phyxminiproblems-problemphyxmini0903-lengthquantity, def:physics:phyx-mini-0903:phyxminiproblems-problemphyxmini0903-lengthinmeters}
  Centimetre readout used by the three side labels in the image
\end{definition}
\begin{proof}
  This is the declared model definition of length In Centimeters.
\end{proof}
\begin{definition}[electric Field Vector In Newtons Per Coulomb]
  \label{def:physics:phyx-mini-0903:phyxminiproblems-problemphyxmini0903-electricfieldvectorinnewtonspercoulomb}
  \lean{PhyXMiniProblems.ProblemPhyXMini0903.electricFieldVectorInNewtonsPerCoulomb}
  \uses{def:physics:phyx-mini-0903:phyxminiproblems-problemphyxmini0903-planevector, def:physics:phyx-mini-0903:phyxminiproblems-problemphyxmini0903-planeelectricfieldquantity}
  Coherent-SI readout of an electric-field vector, in newtons per coulomb
\end{definition}
\begin{proof}
  This is the declared model definition of electric Field Vector In Newtons Per Coulomb.
\end{proof}
\begin{definition}[electric Field Magnitude In Newtons Per Coulomb]
  \label{def:physics:phyx-mini-0903:phyxminiproblems-problemphyxmini0903-electricfieldmagnitudeinnewtonspercoulomb}
  \lean{PhyXMiniProblems.ProblemPhyXMini0903.electricFieldMagnitudeInNewtonsPerCoulomb}
  \uses{def:physics:phyx-mini-0903:phyxminiproblems-problemphyxmini0903-planeelectricfieldquantity, def:physics:phyx-mini-0903:phyxminiproblems-problemphyxmini0903-electricfieldvectorinnewtonspercoulomb}
  Magnitude of the coherent-SI electric-field readout, in newtons per coulomb
\end{definition}
\begin{proof}
  This is the declared model definition of electric Field Magnitude In Newtons Per Coulomb.
\end{proof}
\begin{definition}[Observation Point]
  \label{def:physics:phyx-mini-0903:phyxminiproblems-problemphyxmini0903-observationpoint}
  \lean{PhyXMiniProblems.ProblemPhyXMini0903.ObservationPoint}
  The two observation points printed at the lower vertices
\end{definition}
\begin{proof}
  This is the declared model definition of Observation Point.
\end{proof}
\begin{definition}[Diagram Vertex]
  \label{def:physics:phyx-mini-0903:phyxminiproblems-problemphyxmini0903-diagramvertex}
  \lean{PhyXMiniProblems.ProblemPhyXMini0903.DiagramVertex}
  The three vertices of the triangular diagram
\end{definition}
\begin{proof}
  This is the declared model definition of Diagram Vertex.
\end{proof}
\begin{definition}[Diagram Edge]
  \label{def:physics:phyx-mini-0903:phyxminiproblems-problemphyxmini0903-diagramedge}
  \lean{PhyXMiniProblems.ProblemPhyXMini0903.DiagramEdge}
  The three dashed sides of the triangle
\end{definition}
\begin{proof}
  This is the declared model definition of Diagram Edge.
\end{proof}
\begin{definition}[endpoints]
  \label{def:physics:phyx-mini-0903:phyxminiproblems-problemphyxmini0903-diagramedge-endpoints}
  \lean{PhyXMiniProblems.ProblemPhyXMini0903.DiagramEdge.endpoints}
  \uses{def:physics:phyx-mini-0903:phyxminiproblems-problemphyxmini0903-diagramvertex, def:physics:phyx-mini-0903:phyxminiproblems-problemphyxmini0903-diagramedge}
  Endpoints assigned to each named side of the triangle
\end{definition}
\begin{proof}
  This is the declared model definition of endpoints.
\end{proof}
\begin{definition}[vertex]
  \label{def:physics:phyx-mini-0903:phyxminiproblems-problemphyxmini0903-observationpoint-vertex}
  \lean{PhyXMiniProblems.ProblemPhyXMini0903.ObservationPoint.vertex}
  \uses{def:physics:phyx-mini-0903:phyxminiproblems-problemphyxmini0903-observationpoint, def:physics:phyx-mini-0903:phyxminiproblems-problemphyxmini0903-diagramvertex}
  Vertex corresponding to a printed observation-point label
\end{definition}
\begin{proof}
  This is the declared model definition of vertex.
\end{proof}
\begin{definition}[expected Point Label]
  \label{def:physics:phyx-mini-0903:phyxminiproblems-problemphyxmini0903-expectedpointlabel}
  \lean{PhyXMiniProblems.ProblemPhyXMini0903.expectedPointLabel}
  \uses{def:physics:phyx-mini-0903:phyxminiproblems-problemphyxmini0903-observationpoint}
  Text printed beside each lower observation point
\end{definition}
\begin{proof}
  This is the declared model definition of expected Point Label.
\end{proof}
\begin{definition}[Figure Charge Sign]
  \label{def:physics:phyx-mini-0903:phyxminiproblems-problemphyxmini0903-figurechargesign}
  \lean{PhyXMiniProblems.ProblemPhyXMini0903.FigureChargeSign}
  The sign glyph drawn inside the source circle
\end{definition}
\begin{proof}
  This is the declared model definition of Figure Charge Sign.
\end{proof}
\begin{definition}[Charge Source Kind]
  \label{def:physics:phyx-mini-0903:phyxminiproblems-problemphyxmini0903-chargesourcekind}
  \lean{PhyXMiniProblems.ProblemPhyXMini0903.ChargeSourceKind}
  Physical classification of the displayed charge source
\end{definition}
\begin{proof}
  This is the declared model definition of Charge Source Kind.
\end{proof}
\begin{definition}[Electric Field Triangle Figure]
  \label{def:physics:phyx-mini-0903:phyxminiproblems-problemphyxmini0903-electricfieldtrianglefigure}
  \lean{PhyXMiniProblems.ProblemPhyXMini0903.ElectricFieldTriangleFigure}
  \uses{def:physics:phyx-mini-0903:phyxminiproblems-problemphyxmini0903-observationpoint, def:physics:phyx-mini-0903:phyxminiproblems-problemphyxmini0903-diagramedge, def:physics:phyx-mini-0903:phyxminiproblems-problemphyxmini0903-figurechargesign}
  Literal and idealized presentation data transcribed from image 903.png
\end{definition}
\begin{proof}
  This is the declared model definition of Electric Field Triangle Figure.
\end{proof}
\begin{definition}[Point Charge Triangle Setup]
  \label{def:physics:phyx-mini-0903:phyxminiproblems-problemphyxmini0903-pointchargetrianglesetup}
  \lean{PhyXMiniProblems.ProblemPhyXMini0903.PointChargeTriangleSetup}
  \uses{def:physics:phyx-mini-0903:phyxminiproblems-problemphyxmini0903-signedchargequantity, def:physics:phyx-mini-0903:phyxminiproblems-problemphyxmini0903-planepositionquantity, def:physics:phyx-mini-0903:phyxminiproblems-problemphyxmini0903-lengthquantity, def:physics:phyx-mini-0903:phyxminiproblems-problemphyxmini0903-planeelectricfieldquantity, def:physics:phyx-mini-0903:phyxminiproblems-problemphyxmini0903-observationpoint, def:physics:phyx-mini-0903:phyxminiproblems-problemphyxmini0903-diagramvertex, def:physics:phyx-mini-0903:phyxminiproblems-problemphyxmini0903-chargesourcekind, def:physics:phyx-mini-0903:phyxminiproblems-problemphyxmini0903-electricfieldtrianglefigure}
  The independent physical quantities represented by the image. In particular, electricFieldAt is an observable and is not defined from an answer choice or from the requested 1.8 * 10\textasciicircum{}5 N/C value
\end{definition}
\begin{proof}
  This is the declared model definition of Point Charge Triangle Setup.
\end{proof}
\begin{definition}[displacement From Source In Meters]
  \label{def:physics:phyx-mini-0903:phyxminiproblems-problemphyxmini0903-displacementfromsourceinmeters}
  \lean{PhyXMiniProblems.ProblemPhyXMini0903.displacementFromSourceInMeters}
  \uses{def:physics:phyx-mini-0903:phyxminiproblems-problemphyxmini0903-planevector, def:physics:phyx-mini-0903:phyxminiproblems-problemphyxmini0903-positioninmeters, def:physics:phyx-mini-0903:phyxminiproblems-problemphyxmini0903-observationpoint, def:physics:phyx-mini-0903:phyxminiproblems-problemphyxmini0903-pointchargetrianglesetup}
  Metre displacement vector from the source to an observation point
\end{definition}
\begin{proof}
  This is the declared model definition of displacement From Source In Meters.
\end{proof}
\begin{definition}[edge Length In Meters]
  \label{def:physics:phyx-mini-0903:phyxminiproblems-problemphyxmini0903-edgelengthinmeters}
  \lean{PhyXMiniProblems.ProblemPhyXMini0903.edgeLengthInMeters}
  \uses{def:physics:phyx-mini-0903:phyxminiproblems-problemphyxmini0903-positioninmeters, def:physics:phyx-mini-0903:phyxminiproblems-problemphyxmini0903-diagramedge, def:physics:phyx-mini-0903:phyxminiproblems-problemphyxmini0903-pointchargetrianglesetup}
  Distance in metres between the endpoints of a displayed triangle edge
\end{definition}
\begin{proof}
  This is the declared model definition of edge Length In Meters.
\end{proof}
\begin{definition}[edge Length In Centimeters]
  \label{def:physics:phyx-mini-0903:phyxminiproblems-problemphyxmini0903-edgelengthincentimeters}
  \lean{PhyXMiniProblems.ProblemPhyXMini0903.edgeLengthInCentimeters}
  \uses{def:physics:phyx-mini-0903:phyxminiproblems-problemphyxmini0903-diagramedge, def:physics:phyx-mini-0903:phyxminiproblems-problemphyxmini0903-pointchargetrianglesetup, def:physics:phyx-mini-0903:phyxminiproblems-problemphyxmini0903-edgelengthinmeters}
  Distance in centimetres between the endpoints of a displayed edge
\end{definition}
\begin{proof}
  This is the declared model definition of edge Length In Centimeters.
\end{proof}
\begin{definition}[Matches Single Point Charge Scenario]
  \label{def:physics:phyx-mini-0903:phyxminiproblems-problemphyxmini0903-matchessinglepointchargescenario}
  \lean{PhyXMiniProblems.ProblemPhyXMini0903.MatchesSinglePointChargeScenario}
  \uses{def:physics:phyx-mini-0903:phyxminiproblems-problemphyxmini0903-pointchargetrianglesetup}
  The displayed upper object has the point-charge role used by Coulomb's law
\end{definition}
\begin{proof}
  This is the declared model definition of Matches Single Point Charge Scenario.
\end{proof}
\begin{definition}[Matches Supplied Electric Field Triangle Figure]
  \label{def:physics:phyx-mini-0903:phyxminiproblems-problemphyxmini0903-matchessuppliedelectricfieldtrianglefigure}
  \lean{PhyXMiniProblems.ProblemPhyXMini0903.MatchesSuppliedElectricFieldTriangleFigure}
  \uses{def:physics:phyx-mini-0903:phyxminiproblems-problemphyxmini0903-chargeinnanocoulombs, def:physics:phyx-mini-0903:phyxminiproblems-problemphyxmini0903-lengthinmeters, def:physics:phyx-mini-0903:phyxminiproblems-problemphyxmini0903-lengthincentimeters, def:physics:phyx-mini-0903:phyxminiproblems-problemphyxmini0903-expectedpointlabel, def:physics:phyx-mini-0903:phyxminiproblems-problemphyxmini0903-pointchargetrianglesetup, def:physics:phyx-mini-0903:phyxminiproblems-problemphyxmini0903-edgelengthinmeters, def:physics:phyx-mini-0903:phyxminiproblems-problemphyxmini0903-edgelengthincentimeters}
  The labels, dashed geometry, and charge read directly from the primary image, together with their association to the independent physical quantities. No electric-field value appears in this evidence structure
\end{definition}
\begin{proof}
  This is the declared model definition of Matches Supplied Electric Field Triangle Figure.
\end{proof}
\begin{definition}[Uses School Coulomb Constant]
  \label{def:physics:phyx-mini-0903:phyxminiproblems-problemphyxmini0903-usesschoolcoulombconstant}
  \lean{PhyXMiniProblems.ProblemPhyXMini0903.UsesSchoolCoulombConstant}
  \uses{def:physics:phyx-mini-0903:phyxminiproblems-problemphyxmini0903-pointchargetrianglesetup}
  The rounded coherent-SI value of Coulomb's constant used in the elementary multiple-choice calculation. The scalar itself is Physlib's Electromagnetism.EMSystem.coulombConstant
\end{definition}
\begin{proof}
  This is the declared model definition of Uses School Coulomb Constant.
\end{proof}
\begin{definition}[Has Physical Point Charge Triangle Parameters]
  \label{def:physics:phyx-mini-0903:phyxminiproblems-problemphyxmini0903-hasphysicalpointchargetriangleparameters}
  \lean{PhyXMiniProblems.ProblemPhyXMini0903.HasPhysicalPointChargeTriangleParameters}
  \uses{def:physics:phyx-mini-0903:phyxminiproblems-problemphyxmini0903-chargeincoulombs, def:physics:phyx-mini-0903:phyxminiproblems-problemphyxmini0903-lengthinmeters, def:physics:phyx-mini-0903:phyxminiproblems-problemphyxmini0903-observationpoint, def:physics:phyx-mini-0903:phyxminiproblems-problemphyxmini0903-diagramedge, def:physics:phyx-mini-0903:phyxminiproblems-problemphyxmini0903-pointchargetrianglesetup, def:physics:phyx-mini-0903:phyxminiproblems-problemphyxmini0903-edgelengthinmeters}
  Positivity and source-separation conditions for the intended geometry
\end{definition}
\begin{proof}
  This is the declared model definition of Has Physical Point Charge Triangle Parameters.
\end{proof}
\begin{definition}[Satisfies Planar Point Charge Coulomb Law]
  \label{def:physics:phyx-mini-0903:phyxminiproblems-problemphyxmini0903-satisfiesplanarpointchargecoulomblaw}
  \lean{PhyXMiniProblems.ProblemPhyXMini0903.SatisfiesPlanarPointChargeCoulombLaw}
  \uses{def:physics:phyx-mini-0903:phyxminiproblems-problemphyxmini0903-chargeincoulombs, def:physics:phyx-mini-0903:phyxminiproblems-problemphyxmini0903-electricfieldvectorinnewtonspercoulomb, def:physics:phyx-mini-0903:phyxminiproblems-problemphyxmini0903-pointchargetrianglesetup, def:physics:phyx-mini-0903:phyxminiproblems-problemphyxmini0903-displacementfromsourceinmeters}
  For a source charge q at r₀, the static vector field at r ≠ r₀ is E(r) = k q (r - r₀) / ‖r - r₀‖³. This is a general governing law at both labelled observation points and does not contain the requested numerical magnitude
\end{definition}
\begin{proof}
  This is the declared model definition of Satisfies Planar Point Charge Coulomb Law.
\end{proof}
\begin{lemma}[point1 source distance readout]
  \label{lem:physics:phyx-mini-0903:phyxminiproblems-problemphyxmini0903-point1-source-distance-readout}
  \lean{PhyXMiniProblems.ProblemPhyXMini0903.point1_source_distance_readout}
  \uses{def:physics:phyx-mini-0903:phyxminiproblems-problemphyxmini0903-pointchargetrianglesetup, def:physics:phyx-mini-0903:phyxminiproblems-problemphyxmini0903-edgelengthincentimeters, def:physics:phyx-mini-0903:phyxminiproblems-problemphyxmini0903-matchessuppliedelectricfieldtrianglefigure}
  The figure places point 1 exactly 1.0 cm from the source charge
\end{lemma}
\begin{proof}
  The conclusion follows from the governing relations and figure readouts named in the statement of point1 source distance readout.
\end{proof}
\begin{lemma}[source charge readout]
  \label{lem:physics:phyx-mini-0903:phyxminiproblems-problemphyxmini0903-source-charge-readout}
  \lean{PhyXMiniProblems.ProblemPhyXMini0903.source_charge_readout}
  \uses{def:physics:phyx-mini-0903:phyxminiproblems-problemphyxmini0903-chargeinnanocoulombs, def:physics:phyx-mini-0903:phyxminiproblems-problemphyxmini0903-pointchargetrianglesetup, def:physics:phyx-mini-0903:phyxminiproblems-problemphyxmini0903-matchessuppliedelectricfieldtrianglefigure}
  The physical source-charge readout is the displayed -2.0 nC
\end{lemma}
\begin{proof}
  The conclusion follows from the governing relations and figure readouts named in the statement of source charge readout.
\end{proof}
\begin{lemma}[field At Point1 eq coulomb Vector]
  \label{lem:physics:phyx-mini-0903:phyxminiproblems-problemphyxmini0903-fieldatpoint1-eq-coulombvector}
  \lean{PhyXMiniProblems.ProblemPhyXMini0903.fieldAtPoint1_eq_coulombVector}
  \uses{def:physics:phyx-mini-0903:phyxminiproblems-problemphyxmini0903-chargeincoulombs, def:physics:phyx-mini-0903:phyxminiproblems-problemphyxmini0903-electricfieldvectorinnewtonspercoulomb, def:physics:phyx-mini-0903:phyxminiproblems-problemphyxmini0903-pointchargetrianglesetup, def:physics:phyx-mini-0903:phyxminiproblems-problemphyxmini0903-displacementfromsourceinmeters, def:physics:phyx-mini-0903:phyxminiproblems-problemphyxmini0903-hasphysicalpointchargetriangleparameters, def:physics:phyx-mini-0903:phyxminiproblems-problemphyxmini0903-satisfiesplanarpointchargecoulomblaw}
  Specializing the governing law at point 1 gives the vector expression from which the requested magnitude is later calculated
\end{lemma}
\begin{proof}
  The conclusion follows from the governing relations and figure readouts named in the statement of field At Point1 eq coulomb Vector.
\end{proof}
\begin{definition}[Answer Choice]
  \label{def:physics:phyx-mini-0903:phyxminiproblems-problemphyxmini0903-answerchoice}
  \lean{PhyXMiniProblems.ProblemPhyXMini0903.AnswerChoice}
  Labels attached to the four displayed answer choices
\end{definition}
\begin{proof}
  This is the declared model definition of Answer Choice.
\end{proof}
\begin{definition}[field Magnitude In Newtons Per Coulomb]
  \label{def:physics:phyx-mini-0903:phyxminiproblems-problemphyxmini0903-answerchoice-fieldmagnitudeinnewtonspercoulomb}
  \lean{PhyXMiniProblems.ProblemPhyXMini0903.AnswerChoice.fieldMagnitudeInNewtonsPerCoulomb}
  \uses{def:physics:phyx-mini-0903:phyxminiproblems-problemphyxmini0903-answerchoice}
  Field magnitude in N/C printed beside each answer label
\end{definition}
\begin{proof}
  This is the declared model definition of field Magnitude In Newtons Per Coulomb.
\end{proof}
\begin{definition}[recorded Dataset Answer]
  \label{def:physics:phyx-mini-0903:phyxminiproblems-problemphyxmini0903-recordeddatasetanswer}
  \lean{PhyXMiniProblems.ProblemPhyXMini0903.recordedDatasetAnswer}
  \uses{def:physics:phyx-mini-0903:phyxminiproblems-problemphyxmini0903-answerchoice}
  The answer label recorded in the supplied dataset metadata
\end{definition}
\begin{proof}
  This is the declared model definition of recorded Dataset Answer.
\end{proof}
% --- Archon declaration topology end ---
% --- Archon physics formalization source end ---
